\chapter{\abstractname}

Agent-based transport simulations are the de facto methodology for evaluating urban policy interventions, enabling planners to assess the impact of infrastructure changes on traffic patterns across entire metropolitan areas. However, the computational burden of such simulations---often requiring hours per scenario---severely limits the scope of policy exploration. Graph Neural Networks (GNNs) have been shown to be effective surrogates for these simulations, offering scenario evaluation times reduced from hours to seconds while maintaining reasonable predictive accuracy.

While the predictive performance of GNN surrogates has been studied, their reliability in deployment settings remains underexplored. When a surrogate encounters scenarios that deviate from its training distribution, predictions may become unreliable without any indication to the end user. This limitation poses significant risks for decision support applications where understanding prediction confidence is as critical as the predictions themselves.

In this thesis, we advance the study of uncertainty quantification (UQ) for GNN-based traffic prediction surrogates. We train and evaluate eight model configurations on a Paris road network dataset comprising 10,000 MATSim scenarios; all reported results are based on a 1,000-scenario (10\%) subset unless stated otherwise. The best-performing configuration (Trial~8) achieves $R^2 = 0.5957$, MAE $= 3.96$~vehicles/hour, and RMSE $= 7.12$~vehicles/hour on the test set of 100 graphs (31,635 road segments each), yielding 3,163,500 node-level predictions. A heteroscedastic output extension (Trial~9) was also investigated; training converged to $R^2 = 0.02$ due to variance inflation under the NLL objective and is reported as a negative result.

We systematically compare three UQ approaches. For Trial~8 on its full 100-graph test set, Monte Carlo (MC) Dropout with $S = 30$ stochastic forward passes achieves Spearman uncertainty--error correlation $\rho = 0.4820$. In a separate ensemble comparison experiment (Experiment~A, 5 independently seeded inference runs with corrected weight loading; Section~\ref{sec:ensemble_caveat}), MC Dropout achieves $\rho = 0.4908$ versus ensemble variance $\rho = 0.4370$---a 12.3\% relative improvement at substantially lower computational cost. Conformal prediction at the 95\% level produces well-calibrated coverage (95.01\%) with a quantile of $q = 14.68$~vehicles/hour (primary 50/50 calibration/evaluation split; a 20/80 audit split corroborates at $q = 14.77$), while na\"{i}ve Gaussian intervals using $k = 1.96\sigma$ achieve only 55.6\% empirical coverage---demonstrating that the MC Dropout standard deviation is not a calibrated uncertainty estimate ($k_{95} = 11.34$) without post-hoc correction.

We further demonstrate the practical value of uncertainty-aware prediction: selectively withholding the 50\% most uncertain predictions reduces MAE by 41.2\% (from 3.95 to approximately 2.32~vehicles/hour). Post-hoc temperature scaling ($T = 2.70$) reduces the Expected Calibration Error from 0.265 to 0.048---an 82\% improvement---and adaptive conformal prediction, which normalises nonconformity scores by MC Dropout uncertainty, narrows conditional coverage variation from [62.9\%, 98.6\%] to [90.0\%, 96.2\%] across uncertainty deciles. Analysis via proper scoring rules confirms that the MC Dropout predictive distribution is sharp (CRPS/MAE ratio of 0.857) but underdispersed (PIT Kolmogorov--Smirnov statistic $= 0.245$, excess mass in the tails), and Winkler score comparisons show that conformal intervals substantially outperform na\"{i}ve Gaussian intervals (Winkler 32.3 vs 49.7 at the 90\% level). An $S$-convergence study verifies diminishing returns beyond $S = 30$ forward passes ($\approx 1\%$ improvement in Spearman~$\rho$ from $S = 30$ to $S = 50$). Cross-trial replication on Trial~7 confirms that these findings generalise across model configurations, with conformal prediction restoring near-exact nominal coverage in both cases. We document a significant negative result regarding multi-model ensembling: ensemble experiments with near-identical architectures yield near-zero predictive $R^2$, attributable to a PyTorch Geometric API version incompatibility that caused silent loss of trained parameters during weight loading, rather than a fundamental failure of ensembling as a technique.

Our findings establish MC Dropout combined with conformal calibration as a practical, computationally efficient uncertainty quantification strategy for GNN-based transport surrogates, with direct applicability to uncertainty-aware urban planning decision support.

\vspace{1cm}
\noindent\textbf{Keywords:} Graph Neural Networks, Uncertainty Quantification, Monte Carlo Dropout, Conformal Prediction, Agent-Based Models, Traffic Prediction, Bayesian Deep Learning
